\section{Necessity}

Changes to the constitutional structure are necessary:
\begin{enumerate}[label=\alph*)]
    \item to harden democracy
    \begin{itemize}
        \item against populism
        \item against crime and corruption in elected offices
        \item against democratically elected enemies of democracy
        \item against incompetence of top politicians, above all of the Federal Chancellor
        \item against momentous wrong decisions by the executive and the legislature
    \end{itemize}
    \item to adapt to the proposed new electoral system
\end{enumerate}

The required changes are so far-reaching that it is hardly conceivable to implement them within the normal process of constitutional amendment.
They require legitimation by the people themselves and not only by the people's representatives.
In my opinion, however, it is necessary in any case that the German people give themselves a new constitution in a free decision,
one that draws the lessons from the history of both German states and also takes account of the experience with democracy that has been gained worldwide since the constitution of the Federal Republic was established in 1949.
Within such a constituent process, larger changes to the constitutional structure can also be adopted.
Although I assume that it will actually amount to rewriting the constitution, I will nevertheless only describe the intended changes to the current constitution in the following, so that I do not have to draft a new constitution.
I would also lack the qualifications for that.
After all, I am a software developer and not a political scientist or constitutional lawyer.

\section{Hardening Democracy Against Populism} \label{sec:hardening_democracy}

The struggle against populism must take place on several levels.
It can only truly be defeated by depriving it of its substantive basis.

Apparently, however, this is not so easy, because democracy, too, produces elites and, much worse, a disconnection of the elites from the everyday reality of the non-elites.
The existence of elites as such is rather a good thing: none of us wants to be led by people who are dumber than we are.
Actually, we do not want to be led by people who are just as clever as we are either.
Instead, we want to be led by the best among us.

It becomes problematic when ever larger parts of the population gain the impression that
\begin{enumerate}[label=\alph*)]
    \item the elites live in their own world that no longer has anything to do with reality, or
    \item only pursue their own interests, which run counter to the interests of the \enquote{ordinary citizen}, or
    \item cannot read the times correctly, or
    \item simply have no solution either, or
    \item the elites are not really cleverer than they are, but merely privileged.
\end{enumerate}

The only real way to counter this is by once again bringing people from the bottom to the top.
Unlike monarchy, democracy does allow this without revolution.

However, it does not necessarily make it easy to rise from the bottom to the top either.
With wealthy, educated parents in the right party, one has considerably better chances of advancement than a child from a background far removed from education and politics.

Nevertheless, it is possible for people who know real life to make it to the top.
But it may take a while until someone once again manages to reach a leadership position who is truly a good leader, and not just someone who is a good talker or has the necessary connections.

It is also difficult, though not impossible, to found new parties and let them grow to such a size that they are genuine electable alternatives to the existing ones.

Because democracy permits bloodless upheavals but does not make them easy either, a temporary weakness in leadership, in terms of both personnel and substance, can arise.
This is then easily exploited by populists, people who claim to have simple solutions but in reality make things even worse.

Populists come in two varieties:
\begin{enumerate}[label=\alph*)]
    \item Democratic populists. People with dangerously stupid ideas who nevertheless do not question democracy.
    \item Undemocratic populists. People with dangerously stupid ideas who want to turn democracy into an autocracy.
\end{enumerate}

The way from a) to b) is not far.
When you realize that your stupid ideas do not get through, you can quickly blame democracy and move up from a) to b), as can be seen clearly in the example of Donald Trump's first and second terms.

Whether, and if so how, democracy could be improved further so that positive upheavals would become easier without making negative upheavals easier at the same time, I do not know.
But one can try to make democracy resilient against temporary phases of \enquote{rule of fools}, so as to give the people a little more time to bring forth better leaders and, if necessary, to correct a mistake.

What matters is that the phases are temporary.
The assumption is that the people, when sufficiently desperate about those who have governed so far, can quickly vote a populist party into office, but, like America under Trump's first term, soon realize that they have made a mistake.\footnote{In my assessment, Trump's later re-election is not a correction of the correction, but came about rather through a combination of forgetfulness and despair over the lack of electoral alternatives.}

The democratic structures must then be so stable that one legislative term is not enough to rebuild them, so that in the end there are still free elections at which the party can be voted out again.
It would also be good if the very worst decisions of a populist-dominated Bundestag and a populist Chancellor could still be blocked from outside the Bundestag.

If populist rule extends over several legislative terms, they will probably manage to rebuild democracy according to their ideas over time after all.
It is in the nature of democracy that the people get what they want.
So if they no longer want democracy, they get none.

In a democracy, the people cannot be protected from themselves in the long run.
That is only possible in autocracies.
But they can be protected from the consequences of a lapse.


\section{New Offices and Constitutional Bodies}

\subsection{Constitutional Chamber} \label{subsec:constitutional_chamber}

The Constitutional Chamber is needed to compensate for the deviations from proportionality caused by majority-assured proportional representation.
It is the Bundestag in its composition after the first round of voting and the only institution that can amend the constitution.\footnote{One can discuss whether the Bundesrat should still have to approve all constitutional amendments. For now, I assume that this remains the case.}
If the election was decided after only one round, the Constitutional Chamber is identical to the Bundestag.
Otherwise, it convenes in the composition the Bundestag would have if there had been only one round.
It then also includes candidates who would have been elected after the first round but do not sit in the Bundestag.
It only convenes for its decisions.

For details, see the \mvref{subsec:constitutional-amendments}.

\subsection{Federal Presidential Council} \label{subsec:presidential_council}

It is in the nature of things that constitutional bodies elected by the citizens are particularly susceptible to populism.
Constitutional bodies that are not elected by the citizens can also be infiltrated over time by gradually filling them with enemies of democracy.
But such processes take time.
That is why non-elected constitutional bodies can set limits to populism for a while.

The idea is to let the citizens elect the shaping constitutional bodies (executive and legislature) themselves,
but to have their work controlled by constitutional bodies that are not elected by the citizens.
The controlling constitutional bodies may not themselves shape how the country develops.
After all, they have not been elected.
But they are to check that the elected shapers do their work properly.

We already have one controlling constitutional body.
The Federal Constitutional Court conducts judicial review of legislation.
However, it lacks means of enforcement, and it does not take decisions on matters of state, but only determines what the law is.

I propose strengthening the office of the Federal President in terms of personnel and giving it more authority.
The Federal President should preside over a \enquote{Federal Presidential Council}\footnote{\enquote{State Council} would also be a fitting term. The Federal President as head of state presides over the State Council.
However, the term already carries historical baggage.},
consisting of the President and six other \enquote{Presidential Councillors}.

The Federal President chairs the meetings and decides what goes on the agenda, but otherwise takes decisions on an equal footing with the other members of the Council.

The additional personnel is needed to live up to the increased responsibility for matters of state that is intended for this office.
The Federal Presidential Council is to be able to dissolve the Bundestag or remove the Chancellor in more situations than the Federal President alone can today.
In addition, it should receive veto rights over executive and legislative decisions.

The Federal President retains the representative duties.
The function of head of state, however, is to be shared among seven councillors, including the President.

The Federal Presidential Council is, so to speak, a \enquote{meta-government}.
It does not govern itself, but it installs the government on behalf of the people and can also remove it again if necessary.

\subsubsection{Composition of the Federal Presidential Council} \label{subsubsec:presidential_council_composition}

A controlling constitutional body must be independent of those it is supposed to control.
How to staff institutions that are obliged to be neutral, such as constitutional and supreme courts, offices of the prosecutor general or central banks, with independent people is a long-known problem of democracy\footnote{E.g.\ Bruce Ackerman, \enquote{The New Separation of Powers}, Harvard Law Review 113 (2000), p.~633 ff.; Tarunabh Khaitan, \enquote{Guarantor Institutions}, Asian Journal of Comparative Law 16 (2021), p.~S40 ff.; Mark Tushnet, \enquote{The New Fourth Branch: Institutions for Protecting Constitutional Democracy}, Cambridge University Press 2021.}.
In recent years, there has been growing discussion of how to protect such institutions from being captured or blocked by populists and extremists\footnote{On capture, e.g.\ David Landau, \enquote{Abusive Constitutionalism}, UC Davis Law Review 47 (2013), p.~189 ff.; Kim Lane Scheppele, \enquote{Autocratic Legalism}, University of Chicago Law Review 85 (2018), p.~545 ff.; Steven Levitsky, Daniel Ziblatt, \enquote{How Democracies Die}, Crown 2018.
On blockades, e.g.\ Venice Commission, CDL-PI(2025)023, \url{https://www.coe.int/en/web/venice-commission/-/cdl-pi-2025-023-e}; Gertrude Lübbe-Wolff, \enquote{How to Prevent Blockage of Judicial Appointments}, Verfassungsblog, October~7, 2022, \url{https://verfassungsblog.de/how-to-prevent-blockage-of-judicial-appointments/}; Carolin Lerch, Friedrich Zillessen, \enquote{Parlamentarische Obstruktion}, Verfassungsblog, January~29, 2024, \url{https://verfassungsblog.de/parlamentarische-obstruktion/}.}.
The results of this discussion have already found their way into the 2024 reform of the election of judges to the Federal Constitutional Court.
It can serve as a model, even though it still needs improvement.
How the Federal Presidential Council and the Federal Constitutional Court can be better protected against capture and blockade than before deserves a discussion of its own, which, however, would go beyond the scope of this paper.

\subsection{Federal Office for Opinion Research}\label{subsec:opinion_research_office}

For every government, it is important to know what the people think.
Democratic governments are no exception.
For the proposal made in \cref{subsec:emergency_brake} to take the will of the people into account when deciding on early elections,
a reliable picture of public sentiment is essential.
For such decisions, one cannot rely solely on surveys by private polling institutes.

I therefore propose establishing a Federal Office for Opinion Research that is constantly engaged, without using intelligence-service methods,
in compiling an anonymized picture of public opinion on specific questions and keeping it up to date.
In doing so, it should rely on a larger data basis than the inexpensive representative surveys of private polling institutes.
Their results, however, should definitely also be used as a cross-check.

The Federal Office for Opinion Research is professionally independent, similar to the Federal Statistical Office.
The Office discloses its methodology.

Surveys can be commissioned by the Federal Government, the parliamentary groups of the Bundestag and the Federal Presidential Council.
Government and opposition must have equal access to the Office, not just the government alone.

The commissioning party decides whether and when the results of its survey are published.
There are two reasons for this:
\begin{enumerate}[label=\alph*)]
    \item Government and opposition often want to know what the people think about a question before they take a position or communicate their decision.
    \item Published surveys can trigger an amplification effect, along the lines of: \enquote{If X\% of the population see it that way, then it can't be wrong.}
    If this concerns populist ideas that one wants to counter, one makes life harder for oneself.
    Figures from a state office carry more weight than those from a private institute.
\end{enumerate}

So that the commissioning party cannot cherry-pick the convenient figures, the Office publishes the results itself, and in full, with the wording of the questions, all answers and the methodology.


\section{Succession Rules for the Early Replacement of the Government or Parts of It}\label{sec:succession_rules}

For many reasons, it can happen that the Bundestag and the Federal Government, or only parts of them,
have to be replaced early.
So that the country is never left without a government, there must be replacement rules for all conceivable cases.
The rules currently in force cannot be adopted without modification.
The proposed procedural changes are too fundamental for that.
For example, the Chancellor is now to be elected indirectly by the people and no longer by the Bundestag.
Therefore, the Bundestag cannot simply elect a new Chancellor.
That would only be an option as a temporary interim solution.

In the following, the conceivable cases are therefore gone through one by one:

\subsection{Replacement of the Chancellor Without a Change in the Legislature}\label{subsec:chancellor_replacement_without_legislative_change}
It can happen that the Chancellor has to be replaced early although nothing else has changed in the balance of power
in parliament.

The Chancellor could, e.g., become seriously ill or even die, resign of their own accord
or no longer be allowed to hold office because they have personally committed something so serious
that they may no longer be Chancellor.
In this case, no new parliamentary elections need to be held.
It is sufficient to elect the Chancellor anew.

The party from which the outgoing Chancellor comes must again present three chancellor candidates to the people for election.
Two of the three candidates may, but need not, be the same as in the original federal election.
In a single-party government, the candidate with the most votes automatically becomes Chancellor.
The Federal President must appoint them.

In a coalition government, the procedure from the \mvref{subsec:election-of-lead-candidates-by-the-people} applies.
The Chancellor can be determined by method a) or b) and requires confirmation by the coalition partners
in a secret ballot.

If this confirmation is not given, a new coalition can still be formed analogously to \cref{subsec:coalition_breakup_with_new_coalition}.
The only difference from \cref{subsec:coalition_breakup_with_new_coalition} is that option c) is naturally no longer
available.
Should the old chancellor's party be part of the new coalition, it enters this coalition with the result of the chancellor-candidate by-election.\footnote{The term \enquote{new election of the Chancellor} would not be correct, since whether any of the new chancellor candidates becomes Chancellor at all also depends on the election results of the other coalition partners.}

If forming a new coalition is not wanted or does not succeed, new Bundestag elections must be held.
Until the Federal President appoints a new Chancellor, the Vice Chancellor conducts the business of government on an acting basis.

\runinhead{Note:}
Even if it is clear from the outset that the existing governing coalition will break up over the missing Chancellor,
a chancellor-candidate by-election must still be held if the old chancellor's party takes part in a new coalition.
This applies even if it joins a new coalition not immediately after the old one has broken up,
but, for example, only after one or more further new coalitions have broken up.
The reason is that the party no longer has a chancellor candidate chosen by the people from three candidates.

The Federal President moderates this process.
The President orders a chancellor-candidate by-election or refrains from doing so if a coalition is formed without the old chancellor's party,
or orders new parliamentary elections as soon as it is clear that no successor coalition will be found.
In the case of a single-party government, which precisely does not depend on a coalition to govern, the Federal President
always orders a chancellor-candidate by-election if only the Chancellor drops out.

A single-party government, too, may have to be replaced as a whole, e.g.\ because the party falls out with itself.
This case is to be treated like the breakup of a coalition and to be distinguished from the replacement of the person of the Chancellor described here.

\subsection{Breakup of a Coalition Followed by the Formation of a New Coalition} \label{subsec:coalition_breakup_with_new_coalition}
It often happens that coalitions break up, and sometimes new ones come together afterwards.
Once a new coalition has come together, it must inform the Federal President so that the President orders chancellor-candidate by-elections
in the event that a party cannot present three chancellor candidates elected by the people.
See \cref{subsec:chancellor_replacement_without_legislative_change}.

Once it is ensured that all coalition parties enter the coalition with chancellor candidates elected by the people, the test of a
successful coalition formation is the election of the Chancellor by the coalition in accordance with the \mvref{subsec:election-of-lead-candidates-by-the-people}.
Unlike coalition formations directly after a federal election, however, there should additionally be the option for
the old Chancellor to remain in office if the coalition partners agree on this and the old Chancellor is available
for election.

This applies even if the old Chancellor does not belong to any party of the new coalition and would thus govern without a party of their own
behind them.
The purpose of this rule is to enable maximum stability of the executive.
If one is lucky enough to have a Chancellor who is experienced in office and respected as a person by all
parties, then they should be allowed to continue in office despite party disputes if they themselves
want this and the new coalition partners want it too.

The coalition partners thus have the choice between candidates determined by three different methods:
\begin{enumerate}[label=\alph*)]
  \item The candidate of the strongest coalition party who received the most votes.
  \item The candidate from all coalition parties who received the most votes.
  \item The old Chancellor.
\end{enumerate}

These three methods can yield three different candidates, or only two or one, because, e.g.,
the old chancellor's party is also the strongest party in the new coalition and the old Chancellor is also
the candidate who won the most votes overall in the last federal election.

Should no Chancellor be confirmed in the secret chancellor election (possibly in several rounds), the coalition
has not come about, and the Federal President orders new parliamentary elections as soon as it is clear that no other coalition
will come about and that the incumbent Chancellor will not form a minority government either.

\subsection{Breakup of a Coalition Without the Formation of a New Coalition}
If a coalition supporting the Chancellor breaks up and no new one is formed, this does not necessarily mean that
the Chancellor would have to resign because of it.
The Chancellor can continue to exercise the executive office even without a majority of their own in the legislature, as long as they are not explicitly removed.
If they need the approval of parliament for their plans, they simply have to find suitable majorities for each plan.

If a governing coalition breaks up, no new one is found to replace it and the old Chancellor does not want to lead a minority government
either, the Federal President must order new parliamentary elections.
The old Chancellor must inform the Federal President of this.

\subsection{Replacement of Individual Ministers}
The procedure for replacing federal ministers remains as it is: the Federal President appoints or dismisses them on the proposal of the Chancellor.

\subsection{Replacement of Individual Members of Parliament} \label{subsec:replacement_of_individual_mps}
If individual members of parliament have to be replaced, list candidates are simply succeeded by the next candidates on the list.
For directly elected members, a by-election must be held in the assigned constituency.


\section{Democratic Counterweight} \label{sec:democratic_counterweight}
The proposed electoral procedure gives the government more power.
In return, a democratic counterbalance must be created elsewhere, so that the risk of a government entrenching itself in power is reduced again.

The government has been given more power so that it can fulfil its mandate to shape policy even in challenging times in which much has to be decided.
All the more important is a \enquote{democratic emergency brake} with which a bad government can also be got rid of early if it abuses its power or squanders it through incompetence.
There must be mechanisms to initiate the early replacement of individuals or of the entire government when this is called for.

The challenge is to design the balance of power so that a government has enough power, on the one hand, to push through unpopular decisions as well, and can still be stopped, on the other hand, if it does too much damage.
Even in a democracy, it must be possible to govern against the will of the majority of the population for a time.
If the will of the people is done at all times, there is actually no government at all.
Some measures are simply necessary even though nobody likes them.

The difference from a dictatorship is that governing against the people is only possible for a time and requires personal commitment.
A necessary but unpopular measure can cost you re-election.
That is how it should be, because it reduces the likelihood of abuse of power.

The problem, however, is that some measures are not only unpopular but such serious mistakes that one cannot wait until the end of the legislative term to correct them.

The following sections therefore discuss various replacement scenarios.
While \cref{sec:succession_rules} dealt with succession rules (What happens if, for whatever reason, the government or parts of it have to be replaced early?),
this section deals with how and under what circumstances an early replacement can be initiated.

\subsection{Terminology: Majority Levels} \label{subsec:majority_levels}
Before the actual discussion, a clarification of terms.
Majorities can be defined with different degrees of qualification, e.g.:

\begin{table}[h]
\renewcommand{\arraystretch}{1.4}%
\begin{tabularx}{\linewidth}{| c | l | >{\raggedright\arraybackslash}X | >{\raggedright\arraybackslash}p{3.5cm} |}
  \hline
  \textbf{Level} & \textbf{Name} & \textbf{Meaning} & \textbf{Percentages} \\
  \hline
  1 & Simple majority       & More people are in favor than against. & > 50\% \\
  \hline
  2 & Clear majority        & It is clearly recognizable that more people are in favor than against. Random effects can be ruled out. & approx.~60\% (context-dependent) \\
  \hline
  3 & Two-thirds majority   & Twice as many people are in favor as against. & 66.66\% \\
  \hline
  4 & Overwhelming majority & Many more people are in favor than against. & 80\%\,\dots\,90\% \\
  \hline
  5 & Almost all            & Almost all voters are in favor. Only isolated exceptions against. & 99.9\%\,\dots\,99.999\% (context-dependent) \\
  \hline
  6 & Unanimous             & All voters are in favor. & 100\% \\
  \hline
\end{tabularx}
\caption{Majority levels}
\label{tab:majority_levels}
\end{table}

Wherever I am not sure how strongly qualified the required majority should be, I simply write \enquote{n-th majority} or \enquote{level-n majority}.
This means a majority of level n according to the table above.
In other words, the exact qualification still has to be determined.

The \enquote{majority level} is a parameter one has for balancing the relationship, mentioned in the previous section, between the government's necessary authority and its correctability.
The division into majority levels according to \cref{tab:majority_levels} is not exhaustive.
Levels can be added, or those not needed can be left out.
What should be retained, however, is the principle of describing majorities by their properties and assigning each majority level an integer that increases with the strength of the majority.

Sometimes I will only speak of \enquote{high} and \enquote{low} majorities without specifying exactly which majority levels are meant in that context.
That will then have to be settled in the discussion process that I hope to initiate with this document.

\subsection{Individual Members of Parliament Break the Law} \label{subsec:mps_breaking_the_law}

Individual members of parliament can already be prosecuted today under Article~46 of the Basic Law if the Bundestag consents.
Beyond that, the Bundestag should have the right to suspend members for the duration of the proceedings and, in the case of
proven serious offences, even to withdraw their mandate retroactively, so that the member loses all privileges
associated with a seat in the Bundestag.

Serious offences in this sense are offences
\begin{itemize}
  \item that are deliberately directed against the Federal Republic of Germany, such as espionage for a foreign power, or
  \item offences that, according to the general sense of justice, rule out holding office as a representative of the people, such as violent crimes or sexual offences.
\end{itemize}

Which offences exactly rule out holding office as a representative of the people according to the general sense of justice still has to be determined.
The principle, however, should be clear.

99\% of the population would probably not want to be represented by someone who has committed murder or sexually abused children, however brilliant that person may be otherwise.
On tax evasion and drunk driving, on the other hand, opinions might be divided.

Upon proof of a serious offence, the Bundestag has the right to withdraw the mandate with immediate effect, or even retroactively if the criminal offence already existed throughout the entire
legislative term.
The member concerned has the right to challenge this measure before the Federal Constitutional Court.
Conversely, if there is no majority in the Bundestag for the suspension or withdrawal of the mandate, individual members can bring an action for it before the Federal Constitutional Court, and the
Federal Constitutional Court can grant the action or not.

\subsection{Removal of the Chancellor}

\subsubsection{Conceivable Reasons for Wanting to Remove the Chancellor Early}

Various reasons for wanting to remove the Chancellor are conceivable.
Not all of them should actually lead to removal.

\begin{enumerate}
    \item Incompetence. The Chancellor may simply prove to be incapable.
    \item The Chancellor is no longer physically or mentally equal to the demands of the office.
    \item The Chancellor makes individual decisions that their party/parliamentary group can no longer support.
    \item A coalition breaks up and a new one forms in its place.
    \item The Chancellor makes individual decisions that the n-th majority of the people disapprove of.
    \item The Chancellor makes decisions that lead the country to ruin in the medium to long term.
    \item The Chancellor makes decisions that plunge the country into a catastrophe in the short term. They must be stopped immediately.
    \item The Chancellor breaks the law.
        A distinction must be made between
        \begin{enumerate}
          \item Violations that, because of their seriousness, make it impossible to continue in office. The same considerations apply here as for all members of parliament, only somewhat more strictly, since the Chancellor has much more power and also a greater role-model function than ordinary members of parliament.
            \begin{enumerate}
                \item The Chancellor does not abide by the constitution or does not comply with decisions of the Federal Constitutional Court.
                \item The Chancellor commits high treason, e.g.\ by deliberately betraying state secrets.
                \item The Chancellor commits a crime, such as murder, that according to the general sense of justice disqualifies them from holding the office of Chancellor.
            \end{enumerate}
          \item Violations that can be punished during the term of office without affecting it, e.g.\ if only a fine is due.
          \item Violations that are not so serious that the term of office could not be continued, but that can only be punished after it ends. (E.g.\ because a prison sentence is likely that does not have to be enforced immediately.)
        \end{enumerate}

    \item The Chancellor does not break any law, but does something because of which explicitly
        \begin{enumerate}[label=\alph*)]
            \item an n-th majority of the Bundestag wants to remove them from office or
            \item an m-th majority of the people wants to remove them from office.
        \end{enumerate}
    \item The Chancellor does not keep their election promises.
    \item The Chancellor pursues bad policies.
\end{enumerate}

Individual decisions that the parliamentary group or the people do not support are not in themselves a reason for the early removal of the Chancellor.
Only when a qualified majority of the Bundestag or of the people explicitly wants the Chancellor removed from office is that a reason.
Broken election promises and bad policies are not sufficient on their own either; sadly, they are normal in a democracy.
If you have a better Chancellor, you are lucky.
The price for broken election promises and averagely bad policies is paid at the next election.

\Cref{tab:removal_grounds} shows which instrument applies to which reason.

\begin{table}[H]
\renewcommand{\arraystretch}{1.4}%
\begin{tabularx}{\linewidth}{| >{\raggedright\arraybackslash}X | >{\raggedright\arraybackslash}X | c |}
  \hline
  \textbf{Reason} & \textbf{Instrument} & \textbf{Section} \\
  \hline
  Incompetence & Probationary period; chancellor-candidate by-election at the people's request & \ref{subsubsec:probation}, \ref{subsubsubsec:early_elections_by_popular_demand} \\
  \hline
  Physical or mental condition & Removal by the Federal Presidential Council & \ref{subsubsubsec:removal_for_health} \\
  \hline
  Parliamentary group does not support individual decisions & No removal; the Chancellor must seek other majorities or give in & \ref{subsubsec:removal_by_bundestag} \\
  \hline
  Coalition breakup with formation of a new coalition & Succession rule for the formation of a new coalition & \ref{subsec:coalition_breakup_with_new_coalition} \\
  \hline
  People reject individual decisions & No removal; the price is paid at the next election & -- \\
  \hline
  Ruin in the medium to long term & Weak veto; if repeated, removal by the Federal Presidential Council & \ref{subsubsubsec:veto_rights}, \ref{subsubsubsec:removal_by_council} \\
  \hline
  Short-term catastrophe & Veto of the Federal Presidential Council; immediate removal from office by the Bundestag & \ref{subsubsubsec:veto_rights}, \ref{subsubsec:removal_by_bundestag} \\
  \hline
  Serious breach of the law & Removal by court judgment & \ref{subsubsec:removal_by_court} \\
  \hline
  n-th majority of the Bundestag wants removal from office & Removal by the Bundestag & \ref{subsubsec:removal_by_bundestag} \\
  \hline
  m-th majority of the people wants removal from office & Chancellor-candidate by-election at the people's request & \ref{subsubsubsec:early_elections_by_popular_demand} \\
  \hline
  Broken election promises, bad policies & No removal; the price is paid at the next election & -- \\
  \hline
\end{tabularx}
\caption{Reasons for removal and instruments}
\label{tab:removal_grounds}
\end{table}

\subsubsection{Remedy for Incompetence: A Probationary Period for the Chancellor}\label{subsubsec:probation}
A peculiarity of the office of Chancellor is that there is actually no comparable national office that could be used as a reference for a person's suitability for the chancellorship.
The closest office is that of the Minister-President of a federal state.
But even that is not really comparable.
A person can be a popular \enquote{father of the state} in one federal state, but still be completely unsuited to lead the whole country.

The only true way to test a person's abilities as Chancellor is to let them hold the office for a while.
I therefore propose introducing a probationary period for a new Chancellor.
If someone has not yet passed a probationary period as Chancellor, they initially receive the office on probation.
This means that after a period yet to be determined (proposal: half a legislative term), they must be confirmed in office by the people in order to remain Chancellor.
As in any other job, the hurdle for parting ways during the probationary period is considerably lower than afterwards.
Unlike early elections at the people's request (\cref{subsubsubsec:early_elections_by_popular_demand}), the consent of the Federal Presidential Council is therefore not required.

A referendum is held on the question of whether the incumbent Chancellor may remain in office until the end of the legislative term.
If an n-th majority answers \enquote{no}, the Chancellor has failed the probationary period, and the procedure for replacing the Chancellor under \cref{subsec:chancellor_replacement_without_legislative_change} applies.
Otherwise, they have passed the probationary period and never have to govern on probation again.
This also applies should they be voted out in the meantime and later elected again.
The people now know them in office and can make a conscious decision.

For n, I would propose 2, i.e.\ a clear majority.
A simple majority would be too little, because even a new Chancellor has to push through unpopular decisions in the first term.
The majority level should, however, be lower than for early elections at the people's request, so that an unsuitable Chancellor can be got rid of comparatively easily during the probationary period.

One can discuss whether, after failing the probationary period, a renewed candidacy for Chancellor should be possible again after a waiting period or not.

\runinhead{Arguments against:}

\begin{itemize}
  \item Stability of the political system. The people already know the candidate. Why should they do better next time than this time? If they fail again on the next attempt, new elections have to be held again, and so many early elections are an organizational burden and do not do the stability of the system any good.
  \item The people forget quickly. There is a danger that a completely unsuitable candidate, such as Donald Trump, comes to power again and is worse the second time than the first.
\end{itemize}

\runinhead{Arguments for:}
\begin{itemize}
  \item People can actually change. If the waiting period is long enough, say three legislative terms, which in Germany would correspond to 12 years, the candidate may have matured so much in the meantime that they even become a particularly good Chancellor because they have learned from their first experience.
  \item Shortage of candidates. It is hard to find good leaders. Perhaps after a sufficient waiting period, one can try again with an already known candidate.
\end{itemize}

\subsubsection{Removal of the Chancellor by the Bundestag}\label{subsubsec:removal_by_bundestag}
If the Chancellor departs from their party's line on some questions, this is by no means a reason to remove them.
After all, like all other members of parliament, the Chancellor is bound first and foremost to the people and their own conscience and not to their party.

If their party colleagues no longer want to nominate them as a candidate at the next election, that is their good right.
During a legislative term, however, they cannot replace the Chancellor with someone else.

If the parliamentary group falls out with the Chancellor on some questions, that is not yet a reason to remove them either.
After all, the group does not have to approve the Chancellor's plans.
The Chancellor must then see whether they can perhaps get the missing support from parts of the opposition.
If so, they may still be able to push through their decisions after all.
If not, they simply have to give in.
That is a normal democratic process.

Nevertheless, there may be cases in which the parliamentary group not only disagrees with the Chancellor on individual questions, but wants to get rid of the Chancellor.
If it does so in the course of forming a new coalition, the procedure is already laid down in Section~\ref{subsec:coalition_breakup_with_new_coalition}.

Under the Basic Law currently in force, the parliamentary group could at any time replace the Chancellor with another member of the Bundestag by expressing its lack of confidence in the current Chancellor under Article 67 and electing another Chancellor in their place.

Under the new system, replacing people would not be possible unless it came about as a by-product of the formation of a new coalition.
At most, chancellor-candidate by-elections could be triggered.
Article 67 therefore cannot be retained in this form.

Especially in the present turbulent times, in which nothing can be relied on any more, one must assume, however, that there may be situations in which a Chancellor has to be stopped because they are doing too much damage.

\runinhead{Two cases must be distinguished:}
\begin{enumerate}
    \item The Chancellor must be removed from office with immediate effect.
    \item The Chancellor can remain in office in a caretaker capacity until the successor has been elected.
\end{enumerate}

In the first case, the Bundestag can remove the Chancellor from office with a level-n majority and also appoint an interim successor straight away if it should conclude that the Vice Chancellor cannot carry on the business of government either, e.g.\ because they are involved in the same scandal.
In any case, chancellor-candidate by-elections are triggered, which must take place within a period yet to be determined.

In the second case, a level-m majority is sufficient to remove the Chancellor from office and to trigger the chancellor-candidate by-elections.

For n, I would propose 3, i.e.\ a two-thirds majority, since this is after all a very serious intervention in the leadership of the state.
For m, I would choose 2, i.e.\ a clear majority.
But m~=~1, as under the current rules, would also be possible.
Then the parliamentary group could topple the Chancellor on its own, provided it was united.
Only it could not install a successor itself.

\subsubsection{Removal of the Chancellor by Court Judgment}\label{subsubsec:removal_by_court}

The Basic Law apparently presupposes a certain basic morality, at least among the great majority of members of parliament.
Both the executive and the legislature can violate the Basic Law, but they do not reject democracy and the separation of powers as a whole and will therefore voluntarily allow themselves to be corrected by the Federal Constitutional Court in the event of violations.

If a Chancellor wants to play Donald Trump, parliament will stop them.
Unfortunately, however, one has to assume that this precondition might soon no longer be met.

Should an evil charismatic leader like Hitler actually appear again, whom the masses would willingly follow to their deaths, no system, however sophisticated, could stop him.
He would flood every system, because he would have his followers everywhere, in every authority and every branch of government.

Against a \enquote{foolish shepherd}\footnote{Zechariah 11:15--17} in the manner of Donald Trump, however, democracy could perhaps be \enquote{hardened} to such an extent that it would survive until he disappeared from power again.
It would suffer damage from his term in office, but perhaps the irreversible conversion into an autocracy could still be prevented, so that there would still be a time afterwards.

If parliament does not stop a Chancellor who has got out of control because too many of their followers sit in it, a court could still do so.
For that, however, the concept of the immunity of members of parliament (and thus of the Chancellor) would have to be defined more precisely than before.

People can be showered with all kinds of suspicions and lawsuits that later turn out to be unfounded.
But the necessary cooperation in the investigations and the defense alone cost so much time and energy that the members of parliament can no longer attend sufficiently to their parliamentary work.
It may also be that the members have really committed offences, but ones that were not so serious that punishing them would justify an interruption of their parliamentary work.

To prevent this, Article 46 (2) and (3) of the Basic Law requires the approval of the Bundestag for criminal prosecution and any restriction of the personal freedom of a member.
This article would have to be \enquote{beefed up}, and in addition an authority outside the Bundestag would have to weigh up the disruption to the members against the necessary criminal prosecution,
so that prosecution, where appropriate and necessary, can also be carried out if the Bundestag is biased.

\runinhead{The following must be decided:}

\begin{enumerate}
    \item Are individual investigative measures, such as house searches or testimony in court, permissible?
    (Investigations that have no disruptive effect and do not violate Article 47 of the Basic Law are always permissible.)
    In doing so, the public interest in solving the crime must be weighed against the public interest in undisturbed government work.
    At the same time, the expected effectiveness of the investigative measure for solving the crime must be weighed against the impairment it causes to personal life and to the work of government.
    For example, can one really hope to find evidence in the requested house search, or is it clear beforehand that the evidence, if it exists, will not be found in the house of the accused?
    \item May sentences be enforced during the term of office or only afterwards?
    \item May the office continue to be held?
\end{enumerate}

The deliberations on 1 are classic tasks for judges, with the difference that the public interest must be weighed as well.
The necessary authority to approve an investigative measure against the Chancellor contrary to the majority of the Bundestag can only lie with the Federal Constitutional Court, since it is the only authority that can otherwise also correct parliament.
The procedure for initiating investigations against the Chancellor would therefore in principle be the same as for investigations against ordinary members of parliament: the competent prosecuting authority must apply to the Bundestag for the lifting of immunity.
If this is not granted, individual members of parliament can still apply to the Federal Constitutional Court for it to be lifted.
The Federal Constitutional Court must then carry out the deliberations mentioned above and lift the immunity or not.
The Federal Constitutional Court can never initiate a criminal prosecution or lift immunity on its own initiative.
In this respect, it only acts at the request of members of parliament.
It is still open to discussion whether each individual invasive investigative measure should require approval by parliament or the Federal Constitutional Court, or whether there should only be a general lifting of immunity for a specific case.

2 and 3 are related and can be solved by a catalog of penalties to be adopted by the Bundestag no earlier than for the next legislative term.
A fine can certainly be paid during the term of office, but it is clear that the Chancellor cannot conduct the business of office from prison.
So if the prison sentence cannot be postponed, because the offence is so serious that punishing it cannot be delayed, or because there is a risk of cover-up or flight, the Chancellor cannot remain in office.
Conversely, there may be offences that may not carry a prison sentence but are nevertheless incompatible with remaining in office.
Both cases can be covered by a catalog of possible offences and penalties that assigns to offences and/or penalties an office-retention value and, where applicable, a deferral of the sentence.

In its simplest form, the catalog could, for example, stipulate that the imposition of prison sentences that normally could not be suspended would automatically result in immediate removal from office.
The procedure described in \cref{sec:succession_rules} would then automatically apply.

The catalog could, however, also explicitly name exceptions to the rule: for example, that prison sentences for tax evasion can be postponed until after the end of the legislative term, but that standing for office again is not permitted.

How sensible such a rule would be with regard to setting an example and the integrity of the office is another question.
It would, however, be an analogue to the currently existing rule that the office protects against criminal prosecution.

In any case, the constitution should stipulate that the retention in office and the deferral of sentences of members of the government are regulated in a federal law that may not be amended with effect for the current legislative term.

In the event of disagreements about the interpretation of this law, the Federal Constitutional Court decides, as always.

\subsection{Banning Parties and Forfeiture of Democratic Fundamental Rights}
Under Article~18 of the Basic Law, the Federal Constitutional Court can withdraw democratic fundamental rights in the event of abuse, including the right to stand for election.
Under Article~21 of the Basic Law, the Federal Constitutional Court can declare parties unconstitutional.
It has been proposed\footnote{\url{https://www.bundestag.de/resource/blob/658086/50c871fb51049b9edd4153910a0c9398/WD-3-169-19-pdf-data.pdf}, pp.~9--10} to create a new court for these purposes and to leave the Federal Constitutional Court with judicial review of legislation only.
I support this proposal.
It should be possible to withdraw the mandate of members of parliament who have already been elected by court order and to deprive them of eligibility for election if, according to their aims or their behavior, they are recognizably intent on
abolishing democracy or converting it into an order that runs counter to the fundamental principle of freedom underlying the Basic Law.
The wish to change individual provisions of the constitution is absolutely legitimate.
That is precisely why there is the right to amend the constitution with a two-thirds majority.
And one must also be allowed to campaign publicly for this.

The problem arises when democratically elected people want to effectively abolish democracy or convert it into an order that runs counter to the fundamental principle of freedom underlying the Basic Law.
What this \enquote{fundamental principle of freedom} of the Basic Law is still has to be defined more precisely.
At present, it is only laid down in the constitution by the so-called \enquote{eternity clause}.
After extending and refining it if necessary, one could extract it into a separate document, the \enquote{Principles of the Basic Law}, which may never be changed as long as the Federal Republic of Germany exists -- not even with a two-thirds majority.\footnote{Of course, no legal provision can guarantee its own observance. This document of principles, too, can be changed or abolished. But then it would be visible to everyone
that what followed was in essence a new state, even if it were still called \enquote{Federal Republic of Germany}.}

\subsection{Democratic Emergency Brake -- Control of the Executive and the Legislature by the Federal Presidential Council}\label{subsec:emergency_brake}
\subsubsection{Motivation}\label{subsubsec:emergency_brake_motivation}
As long as parliament functions properly, it can easily stop any individual who has got out of control, including the Chancellor. Unfortunately, one has to fear that this might soon no longer be the case. It is enough for a populist party with a leader in the manner of Donald Trump to win an absolute majority or to be awarded one as a result of the iterative procedure, and the Chancellor could already do all sorts of unreasonable things like Donald Trump and the Bundestag could pass unreasonable laws without anyone being able to stop them.

If the majority of the people actually stands behind these things, or at least thinks they have to be accepted because all the other parties have failed, then not much can be done about it from within.
One can only hope that people wake up before it is too late.

Should this awakening come quickly for the majority of the people, however, it is good if there are democratic structures that can force an early dissolution of the Bundestag if necessary.
It makes no sense to dissolve the Bundestag if it is likely that the same populist party would win again in new elections.
And for social cohesion and democracy as a whole, it would be risky to ban a party that the majority of the population wants.
Should it turn out differently, however, and the majority of the people quickly change their minds once they feel the consequences of populist policies first-hand, it would be good if there were an instrument by which a government could also be chased out of office early from outside the Bundestag.

Such an instrument for initiating early elections from outside parliament would sometimes be desirable for other reasons, too.
Even without populism, a government can sometimes be so bad that one does not want to put up with it until the end of the legislative term.

The challenge is that general dissatisfaction with the government's work must be distinguished from the correction of undemocratic developments.
One cannot simply trigger new elections automatically whenever poll ratings are bad.
Otherwise, we might never stop voting.

Even if the people's dissatisfaction with the government stems only from the government's incompetence and not from populism, it may nevertheless be right
to bring the elections forward if the incompetence is objectively great and better alternatives are in sight.

What is needed is an office that is not subject to populism, because it is not elected by the citizens, that does not govern itself, but monitors the government in the executive and the legislature.
It must be able to assess the overall situation of society and take decisions on matters of state for the good of all.
This also includes weighing up whether to grant the people's wish for early new elections or rather not, because it is to be feared
that a new government would not do any better, or that the same people would be elected again anyway.

This is where the Federal Presidential Council proposed in \cref{subsec:presidential_council} comes into play.
It is not only to harden democracy against populism, but to ensure that the state as a whole functions.
In principle, it is to act like a Federal President made up of seven people.

\subsubsection{Rights of the Federal Presidential Council} \label{subsubsec:council_rights}

\subsubsubsection{Early Bundestag Elections or Chancellor-Candidate By-Elections at the People's Request}\label{subsubsubsec:early_elections_by_popular_demand}

The Federal Presidential Council constantly monitors the situation in society and has the Federal Office for Opinion Research (\cref{subsec:opinion_research_office}) keep it informed about the distribution of opinion.
To this end, the Office regularly determines how many citizens currently want early new elections and publishes the result.
If an n-th majority of citizens wants early new elections, the Federal Presidential Council can dissolve the Bundestag and order new elections.
If the dissatisfaction is directed mainly against the Chancellor, it can instead order only a chancellor-candidate by-election.
But it does not have to do so.
It is at its discretion whether or not to yield to the will of the people on this question.

For n, I would propose 3, i.e.\ a two-thirds majority, so that a government can still push through unpopular decisions, which is sometimes unfortunately very necessary.

\runinhead{Right of application}
As long as the n-th majority has not been reached, the Federal Presidential Council cannot order new elections.
Once it has been reached, the Council can act on its own initiative.
In addition, every single member of the Bundestag and every state government can then apply to it to order new elections.

If the Federal Presidential Council rejects an application, it must give public reasons for doing so.
If it relies on survey results or other data, these must be published so that the reasoning can be verified.

\runinhead{Waiting period}
It is to be expected that the opposition will submit an application every time the wish for new elections rises above the n-th majority.
After a decision by the Federal Presidential Council, no new application is therefore admissible for X months, no matter from whom.
Exempt are applications based on a significant change in the situation, such as a marked increase in the wish for new elections or a new event not covered by the previous reasoning.
Applications received before its decision are answered jointly by the Federal Presidential Council.

\runinhead{No candidacy ban}
Unlike dissolution for endangering the welfare of the state (\cref{subsubsubsec:removal_by_council}), there is no candidacy ban in the case of new elections at the people's request.
The members of parliament have done nothing wrong; they have only lost the voters' confidence.

\runinhead{Alternative: obligation from an m-th majority}
One can discuss whether the Federal Presidential Council should be obliged to order new elections from a higher, m-th majority onwards if this majority is stable for Y months.
It could then only suspend them for narrowly defined reasons, such as a state of defense or an acute crisis, and would have to give public reasons for doing so.
Together, n and m determine how strongly the will of the people enters into the decision on early elections: from the n-th majority onwards it may be taken into account, from the m-th onwards it is binding.

\runinhead{Arguments against:}
\begin{itemize}
  \item So far, for good reason, the people have had no way at all to initiate early elections. Even an option that requires the consent of the Federal Presidential Council is a big step.
  \item From the m-th majority onwards, the Federal Presidential Council no longer decides, but merely confirms. Populist campaigns would have a clear goal: to drive the wish for new elections above the m-th majority.
\end{itemize}

\runinhead{Arguments for:}
\begin{itemize}
  \item The obligation protects against the Federal Presidential Council itself permanently blocking the will of the people, whether out of excessive caution or because it has been infiltrated.
\end{itemize}

For the time being, I would dispense with such an obligation.

\subsubsubsection{Removal of the Chancellor Because of Their Physical or Mental Condition}\label{subsubsubsec:removal_for_health}
The Federal Presidential Council can order an examination of the Chancellor's physical and mental condition by several independent experts.
If a majority of them conclude that the Chancellor is no longer physically or mentally equal to the demands of the office, the Federal Presidential Council can
relieve the Chancellor of office.
If they reach this conclusion unanimously, the Federal Presidential Council must relieve the Chancellor of office.

\subsubsubsection{Veto Rights of the Federal Presidential Council}\label{subsubsubsec:veto_rights}
The Federal Presidential Council should be allowed to veto decisions of the Federal Government or the Bundestag that would pose an existential threat to the welfare of the state.
This would be the case, for example,
\begin{enumerate}[label=\alph*)]
    \item if the Bundestag were to agree to enter into an armed conflict that would in all likelihood result in the destruction of Germany, or conversely to decide not to defend the country when it is attacked, because, for ideological or other reasons, it sympathizes with the state in question or its government.\footnote{Of course, the people have the right to dissolve their state if they so wish, and to integrate into another union of states.
      But that requires a referendum and cannot be done by parliamentary resolution.}
    \item It would also be the case if a decision would directly endanger the unity of the state, e.g.\ by leading to the secession of territories or to unrest or even civil wars, or
    \item if a decision aims to undermine democracy, e.g.\ by overriding the separation of powers.
     The abolition of democracy counts as an existential threat to the state even if the continued existence of the German state were secured, since it would then in essence be a different state from the present one.
\end{enumerate}

The veto of the Federal Presidential Council can be challenged before the Federal Constitutional Court, e.g.\ by arguing that the feared existential consequences would not actually occur.
The challenge cannot, of course, have suspensive effect, since the veto is intended precisely for existentially important decisions that cannot be made and then taken back again.

One can consider whether the Federal Presidential Council should additionally receive a further, weaker veto right that can be exercised under less stringent conditions and can be overridden by the Constitutional Chamber with an n-th majority.
The Federal Presidential Council could use its weak veto generally for decisions that indirectly endanger the welfare of the state or the existence of the state, such as leaving NATO or the EU without an alternative plan or imposing economically senseless tariffs on the rest of the world.
(Of course, one can leave NATO and the EU and levy tariffs on any country if this is done in a planned way as part of a sensible risk assessment.)

\subsubsubsection{Removal of the Chancellor or Dissolution of the Bundestag by the Federal Presidential Council}\label{subsubsubsec:removal_by_council}

The Federal Presidential Council can, on its own initiative, both remove the Chancellor and dissolve the Bundestag if it comes to the conclusion that they are seriously endangering the welfare of the state.
If, for example, the Federal Presidential Council has to veto decisions several times, it can be assumed that this Bundestag or this Federal Government is acting against the welfare of the state.
If judgments of the Federal Constitutional Court are deliberately not implemented, with deadlines not being met even after a warning, this would also seriously endanger the welfare of the state and would justify early new elections.
If a pandemic or other epidemic were spreading and the Federal Government simply did nothing about it, that too would seriously endanger the welfare of the state and would be a reason to dismiss the government.
Veto rights do not help in this case, since the government is precisely doing nothing, and so there is nothing to veto.
Nor does the Federal Presidential Council have the right to propose specific measures or to assess the effectiveness of the measures chosen.
Those are tasks of the government.
But simply doing nothing would be a reason to fire the government.

A Chancellor who has been removed is automatically banned and may not stand for election again within the next X years.
The same applies to all members of the governing parliamentary group if the Bundestag is dissolved.
The Federal Presidential Council can explicitly name exceptions to this rule if individual members have shown by their voting behavior that the accusation of endangering the welfare of the state does not apply to them.

The decision to remove or dissolve can be challenged before the Federal Constitutional Court.
The challenge may have suspensive effect.
The Federal Constitutional Court decides on the urgency.
In the main proceedings, it does not have to decide whether the welfare of the state was actually seriously endangered, but whether, given the facts, it could be assumed that it was.
The Federal Presidential Council retains a margin of discretion.
The Federal Constitutional Court ensures, however, that the Federal Presidential Council cannot fire the Bundestag and the Chancellor arbitrarily.
There must be sound reasoning that everyone can follow.

